\documentclass[11pt]{article}
\usepackage[margin=1in]{geometry}
\usepackage{amsmath}
\usepackage{booktabs}
\usepackage{url}
\setlength{\emergencystretch}{2em}
\title{A certified constant-parameter set of the target and a full-horizon target tube}
\date{Prepared October 8, 2026}
\begin{document}
\maketitle

\section{Request}
The user asked how to use the target motion contract (constant tangential
acceleration $A$ and constant sideslip $\beta$) so that the target's error
tube stops growing from frame to frame, and then asked for an implementation
and a test of its effect. Under the contract the whole target path is fixed by
six constants, so the uncertainty is only in what is not yet known about them.
Set membership over the radar history encloses them; the enclosure only
shrinks, so the target set for any fixed future time is nested from sample to
sample. This report records the implementation, its tests and the closed-loop
measurements. It also records what limits the gain.

\section{Design}
\paragraph{Parameters.} At the published sample the six constants are
$\theta=[x;y;\chi;V;A;\kappa]$ in the true ego body frame (the tracker's
frame), $\kappa=\sin\beta/l_r$. The family is the controller's
\texttt{predictTarget}: $p(\tau)=[x;y]+s\,\mathrm{sinc}(\kappa s/2)
[\cos;\sin](\chi+\kappa s/2)$, $s=V\tau+A\tau^2/2$.

\paragraph{Detections.} Each radar detection $j$ of the 2-s history window is
carried into the present body frame by one of two certified transports, the
one with the smaller radius:
\begin{itemize}
\item gyro: the relative heading by trapezoidal gyro integration, radius
  $\mathrm{radar}+\mathrm{gnss}+2\sin(a_j/2)|\rho_j|$, where $a_j$ grows by
  $n_g\,dt+M\,dt^2/4$ per 12.5-ms step ($M$ the friction bound on yaw
  acceleration); this is about 0.024~rad per second of age;
\item heading: the detection's own certified ego heading, radius
  $\mathrm{radar}+\mathrm{gnss}+2\sin(b_j/2)|\rho_j|$, $b_j\approx0.0075$~rad.
\end{itemize}
The present heading error $\varphi$ and the present GNSS error $c$ are shared
by all past detections and are variables of the linear program; their
second-order effect is added to each radius. Each detection gives the
constraint $|p(t_j;\theta)-z_j-\varphi w_j-c|\le r_j+R_j$, with the Euclidean
ball replaced by its circumscribed octagon.

\paragraph{Linearization and slices.} $p$ is linearized; $R_j$ is the Taylor
remainder of \texttt{targetPositionSupport} with the arc length inflated by
its own error, valid over the current domain box. Over the observer's prior
(course about $\pm0.4$~rad, curvature the whole contract) the remainder
swamps every detection: a first version shrank nothing. The course and
curvature ranges are therefore split into slices (0.1~rad by
0.008~m$^{-1}$, at most 64), each linearized about its own values. A slice
whose program is empty cannot explain the detections and is discarded; the
union of the others is kept. A slice whose solver fails numerically (after a
second algorithm) is kept whole. The computation repeats on the shrunken box
while some component shrinks by more than 10\% (at most three passes).

\paragraph{Domain and nesting.} The domain is the observer's published set,
intersected with the contract and with the previous sample's $A$ and $\kappa$
intervals ($V$ widened by the $A$ interval times the elapsed time). The
published \texttt{predictionErrorSet} is intersected with the result. With
\texttt{projectForecast} the published $A$ and $\beta$ are clipped into the
intersected intervals (course unchanged): the projection onto a convex set
that contains the truth cannot move them away from it.

\paragraph{Controller option.} \texttt{collision.targetTubeOrigin} selects
where the target tube of a tightened collision row starts:
\texttt{"holdStart"} (default) restarts the published set at each hold, as
before; \texttt{"currentSample"} propagates it from the present sample over the
whole prediction, so every tightened row encloses the target for the present
information. The tightened rows carry slack in both cases.

\section{Implementation}
\begin{itemize}
\item \path{estimator/nrmmTargetParameterSet.m} (new): the set.
\item \path{estimator/nrmmControllerErrorBounds.m}: intersection, optional
  projection, \texttt{parameterMembershipSet} on the published estimate.
\item \path{estimator/nrmmPositionErrorBound.m},
  \path{estimator/onlineNrmmTrackingRuntime.m}: options, reset on a new track,
  the previous set carried between published samples.
\item \path{config/nrmmTrackingConfig.m}:
  \texttt{runtime.targetParameterSet} (both switches off by default).
\item \path{config/collisionAvoidanceControllerConfig.m},
  \path{controller/solvePredictiveControl.m}: \texttt{targetTubeOrigin}.
\item \path{scripts/runNonlinearPredictiveSafetyValidation.m}: the trace
  records the published set's half-widths.
\end{itemize}

\section{Tests}
\texttt{runtests('tests')}: 595 of 595 passed on the final code. The new class
\texttt{nrmmTargetParameterSetTest} checks, on a synthetic encounter with
bounded radar, GNSS and gyro errors and a slowly varying heading error inside
its radius: containment of the true constants at three samples and a set
narrower than the observer's prior; nesting of $A$ and curvature across
samples; the unavailable cases (no current detection, no transport bound);
and that detections off the constant-parameter family are reported as
inconsistent. A configuration test covers \texttt{targetTubeOrigin}.

\paragraph{First version.} With one linearization over the observer's prior
the remainder reached 1.4--43~m and no component shrank. Course slices alone
gave course $\pm0.16$~rad and speed $\pm1.4$~m/s after 2~s; adding curvature
slices gave course $\pm0.08$~rad, speed $\pm0.7$~m/s and curvature
$\pm0.01$~m$^{-1}$ after 2~s in the same synthetic probe
(\path{scripts/probeParameterSet.m}), with the truth inside at every sample.

\section{Closed-loop campaigns}
The 84 noisy encounters (seven scenarios at 8 and 15~m/s, seeds 20261003 and
1--5, 2000-hold cap, 1-s recovery dwell) ran in four variants, each encounter
in its own MATLAB process, eight in parallel:
\begin{itemize}
\item \emph{baseline}: \texttt{68883d2} (HEAD) unchanged;
\item \emph{V0}: the new code with both options off;
\item \emph{V1}: the set with forecast projection, \texttt{"holdStart"} tube;
\item \emph{V2}: V1 with the \texttt{"currentSample"} tube.
\end{itemize}
\begin{center}
\begin{tabular}{lrrrr}
\toprule
 & baseline & V0 & V1 & V2\\
\midrule
recovered & 55 & 55 & 57 & 31\\
no solved command & 29 & 29 & 27 & 52\\
collision & 0 & 0 & 0 & 1\\
smallest clearance (m) & 0.272 & 0.272 & 0.306 & 0\\
newly recovered / newly failed vs baseline & -- & 0 / 0 & 4 / 2 & 5 / 29\\
\bottomrule
\end{tabular}
\end{center}
Per run: \path{TARGET_PARAMETER_SET_20261008/campaign-runs.csv}; summary:
\path{comparison.txt}.

\paragraph{Baseline.} This is the first complete noisy campaign on
\texttt{68883d2} (the October 6 run stopped after 5 of 84). All 12 braking-lead
runs have no solution at their first frame (solver time limit), as diagnosed
on October 6: the startup rollout never enters the terminal set. The other 17
failures are \texttt{nonlinearAcceptanceFailed} (11) and
\texttt{pcbfNoNumericalResult} (6). The comparison with 82 of 84 on
\texttt{d1c31d6} is now 55 of 84.

\paragraph{V0} reproduces the baseline run for run (outcome, number of holds and
smallest clearance identical in all 84; two failures report the solver time
limit instead of no numerical result, a load-dependent difference). The
differences below are therefore caused by the options.

\paragraph{Width of the published set} (median half-widths over frames, by
time since the first published set):
\begin{center}
\begin{tabular}{lcccc}
\toprule
 & 0--0.5~s & 0.5--1~s & 1--2~s & 2--4~s\\
\midrule
course (rad), V0 / V1 & 0.49 / 0.33 & 0.41 / 0.23 & 0.34 / 0.16 & 0.40 / 0.17\\
speed (m/s), V0 / V1 & 3.81 / 1.91 & 3.23 / 1.35 & 2.75 / 1.05 & 3.48 / 1.15\\
$A$ (m/s$^2$), V0 / V1 & 1.10 / 1.10 & 1.10 / 1.10 & 1.10 / 1.10 & 1.10 / 1.10\\
curvature (m$^{-1}$), V0 / V1 & 0.034 / 0.034 & 0.034 / 0.034 & 0.034 / 0.034 & 0.034 / 0.024\\
\bottomrule
\end{tabular}
\end{center}
The set halves the course and speed widths. It leaves $A$ at the contract
width throughout and curvature until about 2~s, when most encounters are
decided. The set was available in 7317 of 7380 published frames and the
intersection was never empty. The forecast projection changed the published
$A$ or $\beta$ in 78 of 7380 frames.

\paragraph{V1.} Two more encounters recovered, with 4 new recoveries and 2 new
failures; none of the six involved the projection. With the per-hold tube the
target set enters a row only through one hold of growth, so the intersection
mainly lowers the position radius (about 6.5 to 4~cm). The changes are of the
size the October 5 ablations attributed to the closed loop's sensitivity to
small changes, and they are not evidence of a systematic gain.

\paragraph{V2.} Propagating the set over the whole prediction makes the
tightened rows violated by hundreds of metres early in an encounter: the
primary slack sum exceeded 751~m in the top decile of frames, with medians of
1440~m (first 0.5~s) and 332~m (0.5--1~s). Two consequences:
\begin{itemize}
\item In the prefix the tightened row and the nominal collision row of a node
  share one slack. A large robust slack therefore also relaxes the nominal row,
  and the slack sum is dominated by distant nodes that the ego cannot
  influence. In the 15-m/s accelerating head-on of seed 3 (recovered in the
  baseline and in V1) the vehicles touched at 1.07~s.
\item 29 runs that recovered in the baseline failed: 28 without a solved
  command, mostly \texttt{nonlinearAcceptanceFailed}, and the collision.
\end{itemize}
The full-horizon tube is not usable with the present slack structure; it would
need separate, prioritized slacks for robust and nominal rows.

\paragraph{Cost.} On frames where it ran, the set took 0.92~s (median), 2.29~s
(95th percentile) and 7.96~s (maximum) of MATLAB \texttt{linprog} time per
published sample, against the 50-ms controller period. The controller itself
was unchanged (median 0.017~s).

\section{Why the gain is small}
The offline study of the same day assumed radar-only noise (0.04~m) and a
linearization at the truth; it predicted a guaranteed position half-width of
0.2--0.6~m at the impact time with about 0.5~s left. The certified pipeline
cannot reach that:
\begin{itemize}
\item every detection carried into the present frame carries the ego's own
  certified motion error: about 0.0075~rad times the range with the heading
  transport, or 0.024~rad per second of age times the range with the gyro
  transport (the friction bound on yaw acceleration enters the trapezoid
  error). At 30--45~m this is 0.2--0.4~m per detection, five to ten times the
  radar noise;
\item the observer's prior (course $\pm0.4$~rad, speed $\pm3.3$~m/s, $A$ and
  curvature at the contract) makes the linearization remainder large, and
  slicing removes only part of it;
\item the encounters are decided within about 1.5~s of detection, before $A$
  and curvature can be resolved from 2~s of data.
\end{itemize}

\section{Decision needed}
Both options stay off by default; the published estimate and the controller
are unchanged unless they are enabled. The user should decide whether to keep
the implementation in this form. If the target tube is to be made
non-divergent in practice, the measurements point to:
\begin{itemize}
\item a certified relative-motion bound for the ego over short windows (the
  heading and gyro transports dominate the detection radius);
\item separate, prioritized slacks for robust and nominal collision rows
  before any full-horizon tube;
\item a dedicated linear-program solver or fewer slices for real time.
\end{itemize}
The forecast projection touches the published $A$ and $\beta$, which the user
decided on October 5 to take from the tracker; it is off by default and had no
measurable effect here.

\section{Reproduction}
The scripts are in \path{TARGET_PARAMETER_SET_20261008/scripts}. Each campaign
variant runs on a snapshot of the source with \path{solver/} linked:
\begin{verbatim}
campaign.sh <sourceCopy> <outDir>   # with D=<scripts dir>, VARIANT=<MATLAB struct>
run-variants.sh                      # V1 and V2
python3 compare.py out-baseline out-v0-off out-v1-set out-v2-fullfan
\end{verbatim}
\texttt{runOneEncounter} mirrors \path{scripts/runPotentialFieldCampaign.m} for
one encounter and records per-frame published set widths and slack. The raw
outputs were written to
\path{~/.cache/collisionAvoidance/contract_set_impl_20261008/} and are
disposable; the committed folder keeps the per-run table and the summary. The
source hashes of the variant snapshot are in \path{variant-source-sha256.txt}.

\section{Addendum: the estimator stage removed (October 8, 2026)}
At the user's direction the set-membership stage was removed: it is an
estimation mechanism outside the NRMM observer, which the user does not want
(the same reason as the window fit and the parameter stage of October 5). The
published target set again comes from the observer alone, and
\path{estimator/nrmmTargetParameterSet.m}, its test and its wiring in
\path{nrmmControllerErrorBounds}, \path{nrmmPositionErrorBound},
\path{onlineNrmmTrackingRuntime} and \path{nrmmTrackingConfig} are gone; these
files are again identical to \texttt{68883d2}. The user then had the
controller option \texttt{collision.targetTubeOrigin} removed as well
(\texttt{"currentSample"} is not wanted); \path{solvePredictiveControl},
the controller configuration and its test are again identical to
\texttt{68883d2}. Only the per-frame record of the published set's
half-widths in the validation trace remains. The measurements above are kept
as the record of the experiment.

\end{document}
